% =============================================================================
% SECTION III — Factor Selection and Spanning Regressions
% =============================================================================


% ── SEMAPHORE ──────────────────────────────────────────────────────────────────

The benchmark factor models of Section~\ref{sec:benchmarks} are
designed to capture known systematic exposures, but they leave open
the question of whether a broader set of risk factors can better
explain strategy excess returns and absorb the residual alpha.
Starting from the candidate universe of $75$ tradable risk factors
described in Section~\ref{sec:factor_universe}, we proceed in two
steps. Section~\ref{sec:pca} extracts a small number of latent factors
by principal components, following \citet{connor1986performance,
connor1988risk}, and tests whether these diffuse combinations span
strategy excess returns. Section~\ref{sec:aen} turns to a supervised
alternative---best-subset selection \citep{bertsimas2016best}---that
picks the individual factors most relevant for each strategy from the
same candidate set. As in Section~\ref{sec:benchmarks}, each layer is
subjected to the conditional alpha model
(Section~\ref{sec:conditional_alpha}) and to an out-of-sample
re-estimation of the hedge, and
Section~\ref{sec:factor_discussion} consolidates the alpha estimates
across all layers of the paper.


% ── III.A CANDIDATE FACTOR UNIVERSE ───────────────────────────────────────────

\subsection{Candidate Factor Universe}
\label{sec:factor_universe}

We assemble a candidate set of $75$ risk factors from a comprehensive
review of the empirical asset pricing, intermediary, and macro-finance
literatures. Each factor is selected on the basis of a published
theoretical or empirical link to fixed-income arbitrage returns,
intermediary balance-sheet constraints, or broad market risk
conditions. The factors are organised into six economic categories:
credit risk (Panel~A), liquidity and funding (Panel~B), equity
(Panel~C), volatility (Panel~D), macro-financial (Panel~E), and
interest rates and term structure (Panel~F). The complete list with
definitions, data sources, and references is reported in
Table~\ref{tab:factor_list} (Appendix~\ref{app:factor_list}). This
candidate universe is the common input to both the PCA of
Section~\ref{sec:pca} and the best-subset selection of
Section~\ref{sec:aen}.

Every candidate is tradable: each series is the excess return on an
implementable position---an index, a long-short portfolio, a
derivative overlay, or a DV01-neutral curve trade. State variables
enter through traded counterparts rather than as raw proxies; the
intermediary capital factor of \citet{he2017intermediary}, for
example, enters as the excess equity return on the primary-dealer
holding companies (correlation $0.92$ with the underlying capital
ratio), and aggregate liquidity enters through the traded liquidity
portfolio of \citet{pastor2003liquidity}. The restriction matters for
interpretation: with every regressor a realised excess return, the
intercept $\alpha_i$ in the spanning regressions of this section is a
Jensen alpha---the return on an implementable strategy hedged with
implementable positions---rather than the residual of a projection on
non-traded state variables.


% ── III.B PRINCIPAL COMPONENT ANALYSIS ────────────────────────────────────────

\subsection{Principal Component Analysis}
\label{sec:pca}

When the number of candidate factors is large relative to the sample
size, regressing strategy returns on the full set is unreliable: the
strong collinearity of the candidates inflates the variance of the
coefficient estimates and overfits the sample. The standard response
is to summarise the common variation in a small number of principal
components and test whether these diffuse combinations span strategy
returns; estimating latent factors from a large panel and using them
to measure the abnormal return of a managed position is the
performance-measurement framework of \citet{connor1986performance,
connor1988risk}.

We estimate the components from the full estimation sample rather
than from a trailing window. For an in-sample spanning test---as
distinct from a return \emph{forecast}---using the entire sample to
estimate the loadings is both standard and desirable: it delivers the
most precise estimate of the common-factor space, and with the
cross-sectional dimension $N$ and the time dimension $T$ both large it
yields consistent and asymptotically normal second-stage estimates of
$\alpha_i$ and $\boldsymbol{\beta}_i$ \citep{giglio2021asset}. We
standardise each factor to zero mean and unit variance over the
estimation sample and retain the first $K = 8$ components, the number
selected by the \citet{bai2002determining} $IC_{p2}$ criterion (the
less parsimonious $IC_{p1}$ indicates thirteen). The eight components
capture $64.7\%$ of the variance in the balanced factor panel
(Figure~\ref{fig:pca_scree}), and robustness of the alpha to
$K \in \{5, 8, 11, 13\}$ is reported in Appendix~\ref{app:pca_wk}.

The spanning regression takes the form
\[
r_{i,t} \;=\; \alpha_i \;+\; \boldsymbol{\beta}_i'\,
\mathbf{PC}_t \;+\; \varepsilon_{i,t},
\]
where $r_{i,t}$ is the monthly excess return of strategy $i$ and
$\mathbf{PC}_t$ is the $K \times 1$ vector of contemporaneous
principal component scores. Inference uses Newey--West HAC standard
errors \citep{newey1987simple}. Because the regressors are estimated
principal components, the second-stage standard errors are in
principle subject to a generated-regressor adjustment; under
$N, T \to \infty$ with $\sqrt{T}/N \to 0$ the adjustment is
asymptotically negligible \citep{bai2006confidence}, and the
out-of-sample design below, which re-estimates the components on
past-only data, provides the finite-sample check.

Table~\ref{tab:pca_spanning} reports the estimates. The intercept
remains positive and significant at the $1\%$ level for all three
strategies: $+1.44\%$ p.a.\ for BTP~Italia, $+4.78\%$ for CDS--Bond
Basis, and $+2.21\%$ for iTraxx Combined. The adjusted $\bar{R}^2$ is
low throughout, reaching at most $13.4\%$: the diffuse principal
components do not span strategy excess returns.
Table~\ref{tab:pca_loadings} reports, for each retained component, the
candidate factors with the largest absolute loading. The components
admit a clear economic reading: PC1 is dominated by credit and
equity-market factors, PC2 by global rates and bond returns, and PC3
by swap spreads and dealer CDS.

\paragraph{Out-of-Sample Alpha.}

As in Section~\ref{sec:benchmarks}, the full-sample loadings could
reflect an in-sample-optimal combination of the factors rather than a
return an investor could have earned in real time
\citep{moreira2017volatility, cederburg2020performance}. Panel~A of
Table~\ref{tab:pca_oos} addresses this directly. At each evaluation
date we re-estimate the principal components \emph{and} the spanning
betas on a past-only window---expanding (initialised at $60$ months)
as the baseline and a $60$-month rolling window as a robustness
check---and define the out-of-sample alpha as the mean of the hedged
return, tested with the \citet{newey1994automatic} HAC correction.
Because the components are themselves re-estimated within each window,
the design is invariant to the sign and rotation indeterminacy of
principal components. The out-of-sample alpha remains positive and
significant at the $1\%$ level for all three strategies: $+1.61\%$
p.a.\ for BTP~Italia, $+4.55\%$ for CDS--Bond Basis, and $+1.85\%$ for
iTraxx Combined under the expanding window, close to the in-sample
estimates of Table~\ref{tab:pca_spanning}, so the alpha is not an
artefact of full-sample factor estimation.

Because PCA is an unsupervised method---it extracts factors that
maximise variance in the candidate set, not factors that explain
variation in strategy returns---it cannot identify strategy-specific
risk drivers. This motivates the supervised selection of the next
subsection.


% ── III.C BEST-SUBSET SELECTION ───────────────────────────────────────────────

\subsection{Best-Subset Selection}
\label{sec:aen}

\paragraph{Preprocessing.}

The candidate factors are preprocessed for selection. Six candidates
are near-duplicates of another factor in the universe, with absolute
correlation above $0.95$; one member of each pair is removed, leaving
$p = 69$ candidates in the selection design. Each factor is
standardised to unit $L_2$-norm, so that which factors enter is
invariant to scale; point estimates and inference are subsequently
computed on the original units. Even at $p = 69$, the space of
possible factor subsets---$2^{69} \approx 6 \times 10^{20}$---rules
out exhaustive search and calls for a selection method with an
explicit cardinality control.

\paragraph{Selection.}

For each strategy we solve the best-subset problem of
\citet{bertsimas2016best},
\begin{equation}\label{eq:best_subset}
\min_{\alpha,\,\boldsymbol{\beta}}\;
\frac{1}{T}\sum_{t=1}^{T}\bigl(r_t - \alpha -
\boldsymbol{\beta}'\mathbf{f}_t\bigr)^2
\quad \text{subject to} \quad
\|\boldsymbol{\beta}\|_0 \le k,
\end{equation}
computed with the polynomial-time splicing algorithm of
\citet{zhu2020polynomial}. The $\ell_0$ constraint has two properties
that penalised methods lack and that matter in our design. First, a
factor that enters is not shrunk, so among a cluster of correlated
candidates the constraint keeps the most informative member and drains
its near-duplicates, instead of splitting a common signal across the
cluster or shrinking it toward zero as the LASSO does. Second, the
cardinality is controlled directly: we fix $k = 8$ for parsimony and
cross-strategy comparability, rather than selecting the count by an
information criterion or cross-validation, whose objective is nearly
flat---and whose argmin is therefore ill-determined---on
alpha-dominated series with weak factor signal.
Table~\ref{tab:ml_ksens} (Appendix~\ref{app:aen_sensitivity})
re-estimates the selection at $k \in \{5, 8, 10\}$: the selected sets
are broadly nested and the in-sample fit plateaus, so the conclusions
do not hinge on the choice. All-subsets model comparison has
asset-pricing precedent in \citet{barillas2018comparing}.

\paragraph{Post-Selection Inference.}

Given the selected set $\hat{S}_i$, we re-estimate the spanning
regression by unrestricted OLS on the original (un-standardised)
factors,
\begin{equation}\label{eq:post_ols}
r_{i,t} \;=\; \alpha_i \;+\; \boldsymbol{\beta}_i'\,
\mathbf{f}_{\hat{S}_i,t} \;+\; \varepsilon_{i,t},
\end{equation}
with Newey--West HAC standard errors at the
\citet{newey1994automatic} automatic lag. Because the set is selected
on the same data, we complement the analytic $t$-tests with a
stationary block bootstrap \citep{politis1994stationary} of the
post-selection regression ($9{,}999$ replications, expected block
length nine months) and report the percentile $95\%$ confidence
interval for alpha; post-double-selection
\citep{belloni2014inference}, which guards the alpha against
omitted-control bias from the selection step, is one of the
cross-checks in Appendix~\ref{app:aen_selection_matrix}.

\paragraph{Results.}

Table~\ref{tab:aen_stable_ols} reports the post-selection estimates.
Alpha is significant at the $1\%$ level for all three strategies:
$+1.6\%$ p.a.\ for BTP~Italia (bootstrap CI $[0.8, 2.2]$), $+4.8\%$
for CDS--Bond Basis (CI $[2.9, 6.6]$), and $+2.7\%$ for iTraxx
Combined (CI $[1.3, 3.5]$). With eight factors each, the models
deliver adjusted $\bar{R}^2$ of $22\%$, $21\%$, and $18\%$---above
any benchmark framework of Section~\ref{sec:benchmarks} and the PCA
above---yet absorb none of the three alphas.

The selected sets differ across strategies in ways that mirror each
trade's structure. The BTP~Italia set is built around the German
swap-spread curve (the 2-, 5-, and 10-year DV01-neutral swap-spread
trades), traded liquidity, and defensive equity and bond styles; the
significant loadings are the fixed-income defensive factor
($-0.21$, $t = -4.3$) and quality ($+0.04$, $t = 1.9$). The CDS--Bond
set maps onto the intermediary channel: the strategy loads negatively
on the prime-broker CDS factor ($-0.22$, $t = -2.3$)---it pays off
when dealer credit deteriorates---and positively on the default
premium ($+0.11$, $t = 3.2$) and the industrial bond return ($+0.22$,
$t = 3.2$), the natural exposures of a leveraged long-credit position,
with smaller loadings on high yield, emerging-market debt, currency
trend, and the curvature trade. The iTraxx set is drawn from credit-
cycle, volatility, and style factors: value ($-0.03$, $t = -2.2$),
betting-against-beta ($-0.03$, $t = -3.4$), fixed-income momentum
($-0.04$, $t = -2.7$), and the one-month variance-swap return
($-0.99$, $t = -2.5$), with the credit-cycle and implied-volatility
candidates entering but insignificant. No factor is shared by all
three strategies, and any two sets overlap in at most one factor: the
alpha that survives is not compensation for a single common exposure,
and its origin is strategy-specific.

\paragraph{Out-of-Sample Alpha.}

Panel~B of Table~\ref{tab:pca_oos} applies the recursive design of
Section~\ref{sec:pca} to the selected models: the best-subset factor
set is frozen, only the hedge loadings are re-estimated on the
past-only window, and the realised hedged return is tested with the
HAC correction---the real-time abnormal return of the selected model,
with no per-window re-selection. The out-of-sample alpha is positive
and significant at the $1\%$ level for all three strategies under
both the expanding and the rolling scheme, and of the same order as
the in-sample estimates of Table~\ref{tab:aen_stable_ols}.

Two further checks confirm that the selection itself is not the
source of the result. Re-selecting with the adaptive elastic net of
\citet{zou2009adaptive}---a convex, oracle penalised estimator with
data-driven cardinality---and with post-double-selection recovers
substantially overlapping factor sets and the same alpha verdict
(Appendix~\ref{app:aen_selection_matrix}).


% ── III.D CONDITIONAL ALPHA ───────────────────────────────────────────────────

\subsection{Conditional Alpha}
\label{sec:conditional_alpha}

Table~\ref{tab:alpha_regime_pca_aen} decomposes alpha across the three
stress regimes on the iTraxx Main 5-year level---LOW, MEDIUM, and
HIGH, with the same 60 and 100 basis-point cutoffs used throughout the
paper---under both the PCA and the best-subset controls. In every
cell, alpha rises monotonically from the LOW to the HIGH regime. The
CDS--Bond basis shows the sharpest gradient, and BTP~Italia the
familiar pattern of an alpha concentrated away from calm markets:
under the PCA controls it rises from $+0.4\%$ (insignificant) in LOW
to $+3.3\%$ in HIGH, and under the best-subset controls from $+0.5\%$
to $+3.0\%$. For iTraxx Combined the alpha is significant in every
regime and roughly doubles from LOW to HIGH under both specifications.

The \citet{christopherson1998conditioning} conditional model of
Section~\ref{sec:benchmarks}---the full specification of
equation~(\ref{eq:cfg}), with the lagged standardised stress state
conditioning both the alpha and the loadings, and $\alpha_1$ tested by
the moving-block bootstrap---is re-estimated on each factor layer in
Appendices~\ref{app:pca_conditional} and~\ref{app:aen_conditional}.
Under the PCA controls, $\alpha_1$ is significant at the $1\%$ level
for all three strategies. Under the best-subset controls, it remains
large and significant for the CDS--Bond basis ($+2.55\%$ per standard
deviation, $t = 4.27$), while for BTP~Italia and iTraxx Combined the
continuous slope is absorbed once the loadings themselves are allowed
to vary with stress---the regime gradient of
Table~\ref{tab:alpha_regime_pca_aen} survives, but it is carried by
the conditional loadings rather than by a conditional intercept.
Sub-period stability of the alpha estimates is confirmed in
Appendix~\ref{app:pca_subperiod}, and rolling 36-month alpha estimates
with regime shading are plotted in Appendix~\ref{app:aen_rolling}.


% ── III.E ALPHA ACROSS LAYERS ─────────────────────────────────────────────────

\subsection{Alpha Across Layers}
\label{sec:factor_discussion}

Table~\ref{tab:aen_method_comparison} consolidates the answer to the
question this section and the previous one pose. Each cell reports the
alpha and adjusted $\bar{R}^2$ from one layer of controls: the three
benchmark frameworks of Section~\ref{sec:benchmarks}, the principal
components, and the best-subset selection. Moving down the table, the
explanatory power rises---from $\bar{R}^2$ between $0\%$ and $15\%$
under the benchmarks to $18$--$22\%$ under the best-subset
models---while the alpha barely moves: $+1.6$ to $+2.0\%$ p.a.\ for
BTP~Italia, $+4.0$ to $+4.8\%$ for the CDS--Bond basis, and $+1.8$ to
$+2.7\%$ for iTraxx Combined, significant at the $1\%$ level in every
cell. A supervised search over $75$ tradable candidates, free to pick
whichever eight factors best explain each strategy, absorbs no more of
the alpha than the fixed benchmarks do.

Two conclusions follow. First, the excess returns documented in
Section~\ref{sec:performance} are not compensation for systematic risk
exposures, whether those exposures are specified ex ante, extracted as
latent components, or selected by the data. Second, the alpha that
survives is state-dependent, rising when intermediary funding stress
is elevated, under every layer of controls. The natural question is
why such a premium persists---and whether the three mispricings that
generate it are connected through the balance sheets of the
intermediaries that trade them. Section~\ref{sec:interdependencies}
takes up that question.